\begin{claim}{D.2.21}{Proved}{The exact tenth-stage criterion for the canonical summand}
Keep $X=\Sigma K(\FF_2,2)$, $Y=\Sigma^2K(\FF_2,2)$,
$T=\Omega^2Y$, and the factorization
$\overline h:\Sigma^2X'\to Y$, $X'=X/S^3$, of
\textup{(D.2.19y)}. Let $\chi:S^4\to T$ be its double
adjunct on the bottom four-sphere of $X'$; thus
$h_{\FF_2}(\chi)=a^2$. Write
$\widehat\chi:S^6\to Y$ for its double adjunct.

Fix the canonical map $g_2:C_2X\to T$ and the relative
eight-stage extension supplied by \textup{(D.2.19z)}.
There is an extension through relative dimension ten,
allowing the nine-cell fillings to vary, if and only if
\[
 [\widehat\chi,\widehat\chi]\in4\pi_{11}Y.
 \tag{D.2.21a}
\]
Equivalently, the homotopy Browder square of $\chi$ is zero
in $\pi_9T/4\pi_9T$. This criterion concerns the specified
canonical summand and fixed eight-stage choice. It does not
allow arbitrary changes on the eight-cells, and it does not
decide the separate correction's tenth-stage extension.
\end{claim}
\begin{proof}
Use the labelled configuration models for $C_2X$ and $QX$;
the approximation equivalences apply because $X$ is connected
\cite[Construction~2.4; Theorem~6.1]{May1972}.
The spaces of ordered configurations are connected. Their
maps to the corresponding contractible free symmetric-group
spaces therefore admit relative free CW models with cells
of dimension at least two: attach equivariant two-cells to
kill the fundamental group and then cells of higher dimension.
The label filtration consequently has no relative cells of
weight at least three below dimension eleven, since every
positive label of $X$ has dimension at least three.
Weight one is unchanged. Thus the binary relative complex
is the entire ordinary relative cell model through dimension
ten; a stable splitting alone is not being used to identify
these attaching maps.

Write $a_3,b_4,c_5$ for the three bottom label cells of $X$.
Their integral cellular differentials are
$\partial b_4=2a_3$, $\partial a_3=\partial c_5=0$.
For two labels the configuration spaces are the antipodal
circle and contractible antipodal infinite sphere. Choose
one free $\Sigma_2$-cell $e_i$ in each dimension of the
latter. The relative cells have $i\ge2$, and their
configuration boundaries alternate $1+\tau$ in even
degree and $\tau-1$ in odd degree. The tensor permutation
includes the Koszul sign. The complete list through degree
ten is
\[
\begin{array}{c|l}
8&q=e_2(a_3\otimes a_3)\\
9&r=e_3(a_3\otimes a_3),\quad
  t=e_2(a_3\otimes b_4),\quad
  t'=e_2(b_4\otimes a_3)\\
10&u=e_4(a_3\otimes a_3),\quad
  s=e_3(a_3\otimes b_4),\quad
  s'=e_3(b_4\otimes a_3),\\
 &v=e_2(b_4\otimes b_4),\quad
  e_2(a_3\otimes c_5),\quad e_2(c_5\otimes a_3).
\end{array}
 \tag{D.2.21b}
\]
The relative integral differentials are exactly
\[
\begin{gathered}
 \partial r=-2q,\qquad \partial t=-2q,
 \qquad \partial t'=2q,\\
 \partial s=2r-t+t',\qquad
 \partial s'=-2r+t-t',\qquad
 \partial v=2t+2t',
\end{gathered}
 \tag{D.2.21c}
\]
and the other three degree-ten differentials vanish.
These formulas include the internal differential
$\partial b_4=2a_3$; omitting it would give the wrong
divisibility condition. The full saturated kernel in
degree nine has basis $t-r,t'+r$. The nonzero boundary
columns in that basis are $(-1,1),(1,-1),(2,2)$, with
Smith form $\operatorname{diag}(1,4)$. In particular its
relative ninth homology is $\ZZ/4$.

The quotient $X\to X'$ sends every relative cell in
\textup{(D.2.21b)} except $v$ into $C_2X'$: a label in
the collapsed bottom sphere is removed. It supplies the
specified eight-stage map and a choice of the nine-cell
fillings for which every tenth-stage obstruction except
the coefficient on $v$ is zero.

We identify that coefficient homotopically, including its
normalization. On the bottom four-sphere of $X'$, the first
stabilization is $\Omega^2S^6\to\Omega^3S^7$.
The pair $S^6\to\Omega S^7$ is eleven-connected. Its first
relative generator is the twelve-cell of $J_2S^6$.
The latter is obtained from $S^6\times S^6$ by identifying
the two wedge summands with $S^6$, so its attaching map is
exactly $[\iota_6,\iota_6]$. Double looping shifts the
relative dimension to ten. This gives the Browder square
as the boundary of $v$; the integral boundary $2e_1(b_4
\otimes b_4)$ is its usual Hopf-invariant-two normalization,
not an additional factor of two on the Whitehead square.
The next configuration boundary is $\tau-1=0$ on these
even-dimensional labels and does not change this generator.
Thus the obstruction coefficient is
\[
 \theta=\lambda_1(\chi,\chi)\in\pi_9T,
 \qquad \theta\longleftrightarrow
 [\widehat\chi,\widehat\chi]\in\pi_{11}Y.
 \tag{D.2.21d}
\]

Finally keep the eight-stage extension fixed. The possible
changes of the three nine-cell fillings are arbitrary
$R,T_0,T_1\in M=\pi_9T$. Cellular obstruction theory and
\textup{(D.2.21c)} change the coefficients on $s$ and $v$ by
\[
 d_s=2R-T_0+T_1,\qquad d_v=2T_0+2T_1;
 \qquad d_v-2d_s=4(T_0-R).
 \tag{D.2.21e}
\]
The change on $s'$ is $-d_s$ and the other three coefficients
are unchanged. Consequently killing every coefficient
forces $\theta\in4M$. Conversely, if $\theta=4m$, choose
$R=0$ and $T_0=T_1=-m$; this preserves the zero
coefficients and cancels $\theta$. The identity holds for
every abelian group $M$, including its two-torsion, and
proves both directions. The integral chain matrix and
cochain identity are independently checked in
\texttt{code/D2\_canonical\_relative\_chain.py}.
No secondary contribution from changing an eight-cell
filling has been discarded: those fillings were fixed
throughout the assertion.
\end{proof}
